\documentclass{article}
\usepackage[utf8]{inputenc}
\usepackage[T1]{fontenc}
\usepackage{geometry}
\geometry{
    a4paper,
    left=2.5cm,
    right=2.5cm,
    top=2.5cm,
    bottom=2.5cm
}
\usepackage{graphicx}
\usepackage{float}
\usepackage{hyperref}
\usepackage[english]{babel}
\usepackage{amsmath}
\usepackage{amssymb}
\usepackage{array}
\usepackage{booktabs}
\usepackage{caption}
\usepackage{subcaption}
\usepackage{enumitem}
\usepackage{fancyhdr}
\usepackage{titlesec}
\usepackage{url}
\usepackage{textcomp}
\usepackage{xcolor}
\usepackage{listings}
\usepackage{multirow}
\usepackage{longtable}
\usepackage{tabularx}
\usepackage{times}

% Set up page style
\pagestyle{fancy}
\fancyhf{}
\rfoot{\thepage}
\renewcommand{\headrulewidth}{0pt}

% Title formatting
\titleformat{\chapter}[display]{\normalfont\bfseries}{\chaptertitle\thechapter}{0pt}{\Huge}
\titlespacing*{\chapter}{0pt}{50pt}{40pt}

% Document info
\title{Template Document}
\author{Academic Support Agent}
\date{\today}

\begin{document}
\maketitle

% -------- --------
%  .
% ----------------
%	\begin{frame}{ }
%		%\transfade%  
%		
%	\end{frame}

% -------- :  --------
% ----------------
\section{ }
\subsection{ }
\begin{frame}{ }%=0.85
	\begin{itemize}
		\item  :
		\begin{itemize}
			\item  \LaTeX{}  ,  Beamer ,  PPT ;
			\item  ,  \LaTeX{}  ,  ;
		\end{itemize}
		\item  :
		\begin{itemize}
			\item  !
			\item Github: {\color{BSblue}\url{https://github.com/FvNCCR228/SCU_Beamer_Slide-demo}}
			\item Gitee: \color{BSblue}\url{https://gitee.com/NCCR/SCU_Beamer_Slide-demo}
		\end{itemize}
		%			\framebreak
		%		\pagebreak
		\item  :
		\begin{itemize}
			\item  : 522869009@qq.com
		\end{itemize}
	\end{itemize}
	 Windows 11 , Overleaf ( ,  );
\end{frame}

\begin{frame}{ }
	\begin{itemize}
		\item<1-> \LaTeX  :
		\begin{itemize}
			\item<1->  : TeX Live ( \href{https://mirrors.tuna.tsinghua.edu.cn/CTAN/systems/texlive/Images/}{\color{BSblue} } ) TeXstudio VS Code . TeX Live ,  .  Apple IDE ;
			\item<1->  : Overleaf , Slager  ( ).
		\end{itemize}
		\item<2-> \LaTeX  :
		\begin{itemize}
			\item<2->  : Excel2\LaTeX~(\href{https://www.ctan.org/tex-archive/support/excel2latex/}{\color{BSblue}CTAN Excel2\LaTeX});
			\item<2->  : \href{https://www.latexlive.com/}{\color{BSblue}LaTeX }, \href{https://mathpix.com/}{\color{BSblue}
				Mathpix}~\&~\href{https://mathf.itewqq.cn/}{\color{BSblue} LaTeX}.
		\end{itemize}
		\item<3-> \color{JXred}!!  :
		\begin{itemize}
			\item<3-> \color{JXred}!!  UTF-8 ,  XeLaTeX Biber ;
			\item<3->  ,  (Bug );
			\item<3-> \color{JXred}!!  \LaTeX  ,  ".sty" ( ) ,  ( ).
		\end{itemize}
		\item<4-> \color{JXred}  Adobe Acrobat PDF (Ctrl+L ).
	\end{itemize}
\end{frame}

% -------- :  --------
% ----------------
\section{ }
\subsection{ }
\begin{frame}[fragile]{ }
	\begin{multicols}{2}
		\verb|\uline|\hfill  \qquad\uline{ }\\
		\verb|\uuline|\hfill  \qquad\uuline{ }\\
		\verb|\uwave|\hfill  \qquad\uwave{ }\\
		\verb|\sout|\hfill  \qquad\sout{ }\\
		\verb|\xout|\hfill  \qquad\xout{ }\\
		\verb|\dashuline|\hfill  \qquad\dashuline{ }\\
		\verb|\dotuline|\hfill  \qquad\dotuline{ }
	\end{multicols}
\end{frame}

\subsection{ ,  ,  }
\begin{frame}{ }
	\begin{figure}[h]
		\begin{subfigure}{.4\columnwidth}
			\centering
			\includegraphics%
			[width=.3\columnwidth]{stop-rd.pdf}
			\caption{ }
		\end{subfigure}
		\quad
		\begin{subfigure}{.4\columnwidth}
			\centering
			\includegraphics%
			[width=.3\columnwidth]{stop-gn.pdf}
			\caption{ }
		\end{subfigure}
		\caption{ }
	\end{figure}
	\begin{figure}[h]
		\begin{minipage}[t]{.4\columnwidth}
			\centering
			\includegraphics%
			[width=.3\columnwidth]{stop-rd.pdf}
			\caption{ }
			\label{fig:ZhangmenBtdzt}
		\end{minipage}
		\quad
		\begin{minipage}[t]{.4\columnwidth}
			\centering
			\includegraphics%
			[width=.3\columnwidth]{stop-gn.pdf}
			\caption{ }
			\label{fig:ZhangmenWsdzt}
		\end{minipage}
	\end{figure}
\end{frame}

\cprotEnv\begin{frame}
	\frametitle{ }
	 ,  ,  ,  .  .
	\begin{table}[htbp]
		\centering
		\caption{ }
		\label{tab:YixieGfyy}
		\begin{tabular}{rlc}
			\toprule
			  &   &   \\
			\midrule
			  &   &   \\
			  &  ( ) &   \\
			  &   &   \\
			  &   &  ( ) \\
			\bottomrule
		\end{tabular}
	\end{table}
\end{frame}

\subsection{ }
\begin{frame}[fragile]{ }
	\begin{scucode}[]{A welcome program.}{c}{}{codecode}
#include <stdio.h>
int main() {
    printf("Hello world!\n");
    return 0;
}
\end{scucode}
	\begin{scushow}[]{A welcome program.}{c}{}{codeshow}
#include <stdio.h>
int main() {
    printf("Hello world!\n");
    return 0;
}
\end{scushow}
\end{frame}

\subsection{ }
\begin{frame}[fragile,allowframebreaks]{ }
	\begin{scuproof}{}
		 .
	\end{scuproof}
	\begin{scualgorithm}{ Beamer }{ZengmoXBeamerzs}
		\begin{algorithmic}[1]
		\REQUIRE  \LaTeX  ,  
		\ENSURE  
		\STATE  ,  
		\STATE  ,  
		\STATE  ,  
		\RETURN  ,  ,  
		\end{algorithmic}
	\end{scualgorithm}
	\begin{scudefinition}{ }{MalaoJ}
		 ,  ,  ,  ,  .
	\end{scudefinition}
	\begin{sculemma}{ }{JuanwangSlfz}
		 ,  .
	\end{sculemma}
	\begin{scucorollary}{ }{LangrenSdzyx}
		 ,  ,  .
	\end{scucorollary}
	\begin{scuremark}{}{Zhu}
		\vref{coro:LangrenSdzyx},  ,  .
	\end{scuremark}
	\begin{scucondition}{ }{MianshiLrsdtj}
		\vref{coro:LangrenSdzyx},  ,  .
	\end{scucondition}
	\begin{scuconclusion}{ }{AifeiHdbz}
		 :  .
	\end{scuconclusion}
	\begin{scuassumption}{ }{BianzheBhfh}
		 ,  ,  .
	\end{scuassumption}
\end{frame}

\begin{frame}[fragile,allowframebreaks]{ }
	 $f(v) = \dfrac{\mathrm{d}N}{N\,\mathrm{d}v} = 4\pi \Big(\dfrac{\mu}{2\pi kT}\Big)^{3/2} v^2 \mathrm{exp}\Big(-\dfrac{\mu v^2}{2kT}\Big)$.\\[1ex]
	 \[v_p = \sqrt{\dfrac{2kT}{\mu}} = \sqrt{\dfrac{2RT}{M}}\] $R$ , $M = N_A \mu$ .\\[1ex]
		\begin{equation*}
			\bar{v} = \int_0^\infty vf(v)\,\mathrm{d}v = \sqrt{\dfrac{8kT}{\pi\mu}} = \sqrt{\dfrac{8RT}{\pi M}}
		\end{equation*}
	 
		\begin{equation}
		v_{rms} = \Big(\int_0^\infty v^2f(v)\,\mathrm{d}v\Big)^{1/2} = \sqrt{\dfrac{3kT}{\mu}} = \sqrt{\dfrac{3RT}{M}}
		\end{equation}
	$\symbb{R}\symbbit{R}\symcal{R}\symscr{R}\symfrak{R}\symsfup{R}\symsfit{R}\symbfsf{R}\symbfup{R}\symbfit{R}\symbfcal{R}\symbfscr{R}\symbffrak{R}\symbfsfup{R}\symbfsfit{R}$\\[1ex]
	  
		\begin{multline}
			A=\lim_{n\rightarrow\infty}\Delta x\left(a^{2}+\left(a^{2}+2a\Delta x+\left(\Delta x\right)^{2}\right)\right.\label{eq:reset}\\
			+\left(a^{2}+2\cdot2a\Delta x+2^{2}\left(\Delta x\right)^{2}\right)\\
			+\left(a^{2}+2\cdot3a\Delta x+3^{2}\left(\Delta x\right)^{2}\right)\\
			+\ldots\\
			\left.+\left(a^{2}+2\cdot(n-1)a\Delta x+(n-1)^{2}\left(\Delta x\right)^{2}\right)\right)\\
			=\frac{1}{3}\left(b^{3}-a^{3}\right)
		\end{multline}\\[1ex]
	 
		\begin{align}
			E=m&c^2 & &E=mc^2\\
			E=~&mc^2 & E&=mc^2 \notag \\
			E&=mc^2 & E=~&mc^2\\
			&E=mc^2 & E=m&c^2
		\end{align}
	?
		\begin{equation}
			\begin{aligned}
			\dfrac{\mathrm{d}}{\mathrm{d}t} \symbfit{f} &= \dfrac{\mathrm{d}f_x}{\mathrm{d}t} \symbfit{\hat{i}} + \dfrac{\mathrm{d}\symbfit{\hat{i}}}{\mathrm{d}t} f_x + \dfrac{\mathrm{d}f_y}{\mathrm{d}t} \symbfit{\hat{j}} + \dfrac{\mathrm{d}\symbfit{\hat{j}}}{\mathrm{d}t} f_y + \dfrac{\mathrm{d}f_z}{\mathrm{d}t} \symbfit{\hat{k}} + \dfrac{\mathrm{d}\symbfit{\hat{k}}}{\mathrm{d}t} f_z \\
			&= \dfrac{\mathrm{d}f_x}{\mathrm{d}t} \symbfit{\hat{i}} + \dfrac{\mathrm{d}f_y}{\mathrm{d}t} \symbfit{\hat{j}} + \dfrac{\mathrm{d}f_z}{\mathrm{d}t} \symbfit{\hat{k}} + \big[\symbf{\Omega} \times \big( f_x \symbfit{\hat{i}} + f_y \symbfit{\hat{j}} + f_z \symbfit{\hat{k}}\big) \big] \\
			&= \Big(\dfrac{\mathrm{d}\symbfit{f}}{\mathrm{d}t}\Big)_r + \symbf{\Omega} \times \symbfit{f}(t)
			\end{aligned}
		\end{equation}
	!
		\begin{equation}
			\begin{dcases}
			\oint_l \symbfit{H} \cdot \mathrm{d}\symbfit{l} = \iint_S \symbfit{J} \cdot \mathrm{d}\symbfit{S} + \iint_S \dfrac{\partial\symbfit{D}}{\partial t} \cdot \mathrm{d}\symbfit{S} \\
			\oint_l \symbfit{E} \cdot \mathrm{d}\symbfit{l} = - \iint_S \dfrac{\partial\symbfit{B}}{\partial t} \cdot \mathrm{d}\symbfit{S} \\
			\oint_S \symbfit{B} \cdot \mathrm{d}\symbfit{S} = 0 \\
			\oint_S \symbfit{D} \cdot \mathrm{d}\symbfit{S} = \iiint_V \rho \mathrm{d}V
			\end{dcases}
		\end{equation}
	 
		\[
		\begin{bmatrix}
		E & 0 \\
		0 & E \\
		\end{bmatrix}
		\begin{pmatrix}
			E & 0 \\
			0 & E \\
		\end{pmatrix}
		\]
\end{frame}

\subsection{ }
\begin{lrbox}{\mysavebox}
	\begin{scuremark}{ }{JuanwangSlfz}
		 \vref{lemm:JuanwangSlfz}.\par
		\textbf{\color{JXred} :}
		\begin{enumerate}
			\item  ;
			\item  ,  .
		\end{enumerate}
		\textbf{\color{JXred} :}
		\begin{enumerate}
			\item  :  ;\\
			~~ .  ,  :  ,  .  .  
			\begin{itemize}
				\item  
				\item  
				\item  
				\item  
				\item  
				\item  
				\item \dots\dots
			\end{itemize}	
			\item  :  ,  .  ( ,  ) ,  .
		\end{enumerate}
	\end{scuremark}
\end{lrbox}
\begin{frame}[fragile,t,allowframebreaks]{ }
	\framenew{.82}
\end{frame}

\begin{frame}[fragile,label={fra:Fenlan}]{ }
	\vspace{1ex}
	\begin{columns}
		\begin{column}[t]{0.35\linewidth}
			 
			\begin{figure}[h]
				\begin{center}
					\includegraphics[width=.8\columnwidth]{logo_name}
					\\ \\
					\vspace{1ex}
					\includegraphics[width=.4\columnwidth]{Fy}
					\\ 
				\end{center}
			\end{figure}
		\end{column}
		\pause
		\begin{column}[t]{0.4\linewidth}
			 
			\vspace{1ex}
			\begin{itemize}
				\item  
				\begin{enumerate}
					\item  
					\item  
					\item  
				\end{enumerate}
				\item  
				\item  
			\end{itemize}
		\end{column}
		\pause
		\begin{column}[t]{0.25\linewidth}
			 
			\par\vskip1ex
			 \\
			\vskip2ex
			 \\[.5ex]
			 
		\end{column}
	\end{columns}
\end{frame}

\subsection{ }
\begin{frame}[fragile]{ }
	 Beamer ,  .
	\begin{table}[h]
		\centering
		\caption{ }
		\label{tab:JiaochaYymlb}
		\begin{tabular}{lcl}
			\toprule
			  &   &   \\
			\midrule
			\verb|\ref|\verb|{<label>}| &   & \ref{assu:BianzheBhfh} \\
			\verb|\ref*|\verb|{<label>}| &   & \ref*{assu:BianzheBhfh} \\
			\verb|\nameref|\verb|{<label>}| &   & \nameref{assu:BianzheBhfh} \\
			\verb|\vref|\verb|{<label>}| &   & \vref{fra:Fenlan} \\
			\verb|\pageref|\verb|{<label>}| &   & \pageref{assu:BianzheBhfh} \\
			\verb|\vpageref|\verb|{<label>}| &   & \vpageref{assu:BianzheBhfh} \\
			\verb|\cref|\verb|{<label>}| &   & \cref{assu:BianzheBhfh} \\
			\verb|\crefrange|\verb|{<label>}| &   & \crefrange{fig:ZhangmenBtdzt}{fig:ZhangmenWsdzt} \\
			\bottomrule
		\end{tabular}
	\end{table}
\end{frame}

\begin{frame}[fragile]{ }
	\begin{itemize}[<+-| alert@+>]
		\item  \footnote{ .};\\
		\item  \footnotemark.
	\end{itemize}
	\footnotetext{ .}
	 \cite{__2020-1},  \cite{__2016,m_possibilities_2018};\\
	\begin{itemize}
		\item  \footnotemark.\footfullcitetext[3]{_ugc_2018}
		\item  \footfullcite[8]{__2018}.\\
	\end{itemize}
\end{frame}

% -------- :  --------
% ----------------
\section{ }
\subsection{A}
\begin{frame}{Tcolorbox\&Animate}
	\begin{minipage}[c]{0.725\columnwidth}
		\begin{scutcbrb}{ }
			 :  .
		\end{scutcbrb}
		\begin{scutcbyg}{ }
			 :  .
		\end{scutcbyg}
	\end{minipage}~~~~
	\begin{minipage}[c]{0.225\columnwidth}
		\begin{tcolorbox}[enhanced,%
			size=minimal,auto outer arc,%
			width=2.1cm,octogon arc,%
			colback=red,colframe=white,colupper=white,%
			fontupper=\fontsize{7mm}{7mm}\selectfont\bfseries\sffamily,%
			halign=center,valign=center,%
			square,arc is angular,%
			borderline={0.2mm}{-1mm}{red}]
			STOP
		\end{tcolorbox}
	\end{minipage}
%	\begin{center}
%		\begin{animateinline}[loop]{20}%
%			\multiframe{160  }{rx=0+5}{
%				\begin{tikzpicture}[scale=1.5]
%				\pgfmathsetmacro{\x}{4+1*sin(\rx)}	
%				\coordinate  (P) at (\x,0);	
%				\draw [decorate,decoration={coil,segment length= \x pt, pre length=2mm, post length=2mm
%					,amplitude=2mm}] (0,0)--(P);
%				\shade[ball color=black](P)circle (2mm);
%				\fill [pattern = north east lines] (0,0.-0.2) -- ++ (0,0.4) -- ++ (-0.1,0) -- ++(0,-0.4) --cycle
%				(0,-0.2) -- ++(0,-0.1) -- ++ (5.5,0) -- ++ (0,0.1) -- cycle;
%				\draw (0,0.2) -- (0,-0.2) -- (5.5,-0.2);
%				\end{tikzpicture} 	    }
%		\end{animateinline}
%	\end{center}
\end{frame}

%\subsection{B}
%\begin{frame}{Tikz}
%	\begin{minipage}[c]{0.35\linewidth}
%%		\scalebox{20}{}
%		\begin{figure}[h]
%			% Content from tikz-huaji would be included here

%			\vskip-5ex
%			\caption{  -  }
%		\end{figure}
%	\end{minipage}\quad
%	\begin{minipage}[c]{0.6\linewidth}
%		\scalebox{20}{}
%		\begin{figure}[h]
%			% Content from tikz-Dipolar_magnetic_field would be included here

%			\vskip-5ex
%			\caption{Dipolar Magnetic Field - Cyril Langlois}
%		\end{figure}
%	\end{minipage}
%\end{frame}

% -------- :  --------
% ----------------
\section{ }
\subsection{ }
\begin{frame}[allowframebreaks]{ }
	\nocite{*}
	\printbibliography[heading=none]
\end{frame}

% -------- :  --------
% ----------------
\section{ }
\subsection{ }
\begin{frame}{ }
	\begin{center}
		 Beamer, Tcolorbox ,  \\[1ex]
		 , Stack Overflow ,  \\[1ex]
		 ------ :  (Github \href{https://github.com/lxgw/LxgwWenKai/}{\color{BSblue}LxgwWenKai}  ),  : Source Han Sans (Github \href{https://github.com/adobe-fonts/source-han-sans/}{\color{BSblue}source-han-sans}  ),  \\[1ex]
		 Beamer (Github \href{https://github.com/ysx2000/USTCBeamerSX/}{\color{BSblue}USTCBeamerSX}  ),  \\[1ex]
		 Beamer (Github \href{https://github.com/tuna/THU-Beamer-Theme/}{\color{BSblue}THU-Beamer-Theme}  ),  Beamer (Github \href{https://github.com/ustctug/ustcbeamer/}{\color{BSblue}ustcbeamer}  ),  \\[1ex]
		
		 Bug,  ,  \\[1ex]
		 ,  !
	\end{center}
\end{frame}

\begin{frame}{~}
	\centering
	\Huge  
\end{frame}
\end{document}